\chapter{竞品与需求分析}
% 官方要求：结合往年比赛视频、其它学校开源资料、论坛等信息，分析其它学校机器人各项功能的完成度及技术水平，概述值得借鉴之处或者有待改进之处。
\section{华南农业大学——Taurus战队}
\subsection{各项功能的完成度及技术水平}

\subsubsection{机械}
\begin{itemize}
      \item 底盘运动：轮系结构简便、强度高、能满足基本全向移动、爬坡通过角适中、轮组整体对称分布；鲁棒性好，内部摩擦损耗小；陀螺状态稳定，云台晃动小。
      \item 车架设计：采用薄壁铝方管结合碳板进行连接，保证强度同时实现轻量化；防撞设计近似圆形，通过性好且不易在碰撞中被卡住。
      \item Pitch轴设计：采取GM6020电机直连，使云台响应快，没有多余设计；云台转动惯性小。
      \item Yaw轴设计：弹仓跟随Yaw轴转动，转动惯量较大，弹链长度适中，没有明显卡弹现象。
\end{itemize}

\subsubsection{策略}
\begin{itemize}
      \item 路径规划：推测其使用 ros 中的 navigation stack 实现规划。在与联盟赛复杂度类似的场地中能灵活的去到对应的目标点以及避开障碍。
      \item 定位： 将雷达的三维点云信息转换成二维的 laser scan 信息，从而与提前建好的二维地图进行匹配来进行定位。
      \item 决策：能够不断发送目标点。云台能够$\ang{360}$巡航寻找敌人，在看到敌方机器人后能进行自动瞄准。      
\end{itemize}

\subsection{值得借鉴之处}
\begin{itemize}
      \item 决策上采取一边小陀螺一边接近目标点的方式，可以有效规避来自敌方的打击。
      \item 使用雷达的二维信息进行定位，在无高地的环境中表现良好，较为稳定。
   \end{itemize}

\subsection{有待改进之处}
\begin{itemize}
    \item 定位上采用雷达的二维信息定位，在类似分区赛这类存在高地的地形无法适用。
    \item 一边小陀螺一边向目标点虽然能规避地方打击，但也会使哨兵灵活性减少，在较大的场地上可能会来不及支援友方兵种。
\end{itemize}


\section{华南理工大学——松灵华南虎战队}
\subsection{各项功能的完成度及技术水平}

\subsubsection{机械}
\begin{itemize}
      \item 底盘设计：采用舵轮底盘。运动过程中加速度大，响应快。
      \item 弹舱设计：采用两级弹舱。一级弹舱处于底盘，能容纳大量弹丸；二级弹舱固定于yaw轴，能容纳少量弹丸。两级弹舱均有拨弹机构，其中一级弹舱中的拨盘将弹丸拨入二级弹舱，二级弹舱中的拨盘负责将弹丸拨入发射机构中。相较于半上供弹步兵，两级弹舱的设计使得yaw轴惯量更小，整车重心更低。
      \item 采用焊接薄壁铝方管结合碳板进行连接，保证强度同时实现轻量化；防撞设计近似圆形，通过性好且不易在碰撞中被卡住；采用平行四边形连杆的单摇臂独立悬挂，保证轮组相对地面直立，经过起伏路段和下台阶姿态良好。
      \item Pitch轴设计：采取6020电机直连，使云台响应快，没有多余设计；云台转动惯性小；换测速机构方案采用类似于哈尔滨工程大学的双枪转轮机构，占用空间小，响应快。
\end{itemize}

\subsubsection{策略}
\begin{itemize}
      \item 在底盘放置了一个相机用于识别队友装甲板位姿，使得哨兵能够跟随友方兵种。
\end{itemize}

\subsection{值得借鉴之处}
\begin{itemize}
      \item 双级云台的设计能够使云台惯量较半上供云台降低的同时，提高弹舱容量，且不易卡弹。
      \item 使用换测速机构能使哨兵持续高射频射击，且结构简单。
\end{itemize}

\subsection{有待改进之处}
\begin{itemize}
    \item 哨兵采用跟随队友的方案，局限性较大，放弃了更多的战术可能。
    \item 虽然一级弹仓置于底盘，但是仍存在弹仓置于发射机构所在的云台之上，惯量相当于下供弹的方式还是较大，云台响应不够极致。
\end{itemize}